%======================================================================
% CAPÍTULO 16 - SISTEMA DE DISTRIBUCIÓN DE APORTES
%======================================================================
\chapter{Sistema de Distribución de Aportes}

\bigskip

En algunos escenarios, el monto total de garantías calculado excede el presupuesto disponible. Este capítulo describe el algoritmo de distribución proporcional implementado en Sentinel.

\bigskip

\section{Contexto del Problema}

Cuando el presupuesto disponible para garantías es menor que el total calculado, se requiere distribuir el monto aprobado proporcionalmente entre todos los beneficiarios.

\bigskip

El problema crítico es garantizar que la suma de todos los pagos sea exactamente igual al monto aprobado, incluso cuando hay decimales involucrados.

\bigskip

\section{Algoritmo de Distribución Exacta}

El algoritmo implementado garantiza precisión absoluta mediante el uso de aritmética de enteros.

\bigskip

\subsection{Principios del Algoritmo}

\begin{enumerate}
    \item Factor Global Único: Se calcula un factor de proporción único basado en el cociente entre el monto aprobado y la suma total de garantías.
    \item Uso de Centavos (Enteros): Para evitar errores de precisión de punto flotante, todos los cálculos se realizan convirtiendo a centavos (multiplicando por 100) y trabajando con enteros (i64).
    \item Ajuste del Último Registro: El último beneficiario absorbe la diferencia residual para garantizar el cuadre exacto.
\end{enumerate}

\bigskip

\[
\text{Factor Global} = \frac{\text{Monto Aprobado}}{\text{Suma Total de Garantías}}
\]

\bigskip

\subsection{Pseudocódigo del Algoritmo}

\begin{lstlisting}[language=rust]
fn distribuir_garantias(beneficiarios: &mut [Beneficiario], monto_aprobado: f64) {
    // 1. Calcular suma total de garantias
    let suma_garantias: f64 = beneficiarios.iter()
        .map(|b| b.base.garantia_original)
        .sum();

    // 2. Calcular factor global
    let factor_global = monto_aprobado / suma_garantias;

    // 3. Convertir monto aprobado a centavos
    let monto_centavos = (monto_aprobado * 100.0).round() as i64;

    // 4. Distribuir con truncamiento
    let mut acumulado: i64 = 0;
    let n = beneficiarios.len();

    for (i, ben) in beneficiarios.iter_mut().enumerate() {
        let anticipo = ben.base.garantia_original * factor_global;
        let anticipo_centavos = (anticipo * 100.0).round() as i64;

        if i < n - 1 {
            // Truncar (redondeo hacia abajo)
            ben.base.garantia_anticipo = anticipo_centavos as f64 / 100.0;
            acumulado += anticipo_centavos;
        } else {
            // Ultimo: cuadra exacto
            let ajuste = monto_centavos - acumulado;
            ben.base.garantia_anticipo = ajuste as f64 / 100.0;
        }

        ben.base.factor_aplicado = factor_global;
    }
}
\end{lstlisting}

\bigskip

\section{Análisis del Algoritmo}

Teorema de Conservación: La suma de todos los anticipos distribuidos es exactamente igual al monto aprobado.

\bigskip

Demostración: Sea n el número de beneficiarios, M el monto aprobado, gi la garantía original del beneficiario i, y ai el anticipo calculado.

Para i = 1, ..., n-1:
\[
a_i = \left\lfloor \frac{g_i \times M}{\sum g_j} \times 100 \right\rfloor / 100
\]

Para i = n:
\[
a_n = M - \sum_{i=1}^{n-1} a_i
\]

Entonces:
\[
\sum_{i=1}^{n} a_i = \sum_{i=1}^{n-1} a_i + \left( M - \sum_{i=1}^{n-1} a_i \right) = M
\]

\bigskip

\section{Ejemplo Numérico}

\begin{center}
\begin{tabular}{ll}
\toprule
\textbf{Parámetro} & \textbf{Valor} \\
\midrule
Monto Aprobado & 40.000.000,00 Bs \\
Total Garantías Calculadas & 57.049.791,95 Bs \\
Factor Global & 0,70114191 \\
\bottomrule
\end{tabular}
\end{center}

\bigskip

\begin{center}
\begin{tabular}{llllll}
\toprule
\textbf{\#} & \textbf{Cédula} & \textbf{Garantía Original} & \textbf{Cálculo} & \textbf{Anticipo} \\
\midrule
1 & 10002142 & 615,90 & 615,90 × 0,70114191 & 431,83 \\
2 & 10002885 & 611,71 & 611,71 × 0,70114191 & 429,02 \\
3 & 10002920 & 628,93 & 628,93 × 0,70114191 & 440,99 \\
4 & 10009822 & 636,81 & 636,81 × 0,70114191 & 446,51 \\
5 & 10010055 & 593,31 & 593,31 × 0,70114191 & 416,01 \\
... & ... & ... & ... & ... \\
n & ÚLTIMO & X & Ajuste & \textbf{40.000.000,00} \\
\midrule
  & \textbf{TOTAL} & \textbf{57.049.791,95} & & \textbf{40.000.000,00} \\
\bottomrule
\end{tabular}
\end{center}

\bigskip

\section{Verificación}

El algoritmo ha sido verificado mediante pruebas unitarias que confirman la conservación exacta del monto aprobado para diferentes configuraciones de datos.

\bigskip

Ejemplo de prueba unitaria:
\begin{lstlisting}[language=rust]
#[test]
fn test_distribucion_exacta() {
    // Crear registros de prueba
    let mut bases: Vec<Base> = vec![
        create_base(100.0),
        create_base(100.0),
        create_base(100.0),
    ];
    
    // Aprobar solo 150 Bs (la mitad)
    generar_calculos_beneficiarios(&mut bases, 150.0);
    
    // Verificar suma exacta
    let suma: f64 = bases.iter().map(|b| b.garantia_anticipo).sum();
    assert!((suma - 150.0).abs() < 0.01);
}
\end{lstlisting}

\vfill
\newpage
